\documentclass[11pt,a4paper]{article}
\usepackage[margin=2.5cm]{geometry}
\usepackage{amsmath,amssymb}
\usepackage{booktabs}
\usepackage{hyperref}

\newcommand{\bB}{\boldsymbol{B}}
\newcommand{\bD}{\boldsymbol{D}}
\newcommand{\bI}{\boldsymbol{I}}
\newcommand{\bM}{\boldsymbol{M}}
\newcommand{\bN}{\boldsymbol{N}}
\newcommand{\bQ}{\boldsymbol{Q}}
\newcommand{\bR}{\boldsymbol{R}}
\newcommand{\bp}{\boldsymbol{p}}

\begin{document}

\title{Cascade Calculation Methods in \texttt{hydrocal/src/CoolerFractions.cxx}}
\author{AI generated and edited by S. Schippers}
\date{\today}
\maketitle

\section{Problem}

An ion beam is created in hydrogenic excited states $|n,l\rangle$ and
traverses a cooler section of length $L$ before entering a field ioniser.
During the flight of duration $t = L/v_{\rm ion}$ the states undergo
radiative cascade.  We need the \emph{detection probability}
\begin{equation}
  f_{nl} \;=\; \sum_{n',l'} \bigl[\exp(\bM\, t)\bigr]_{n'l',\,nl}\;
  p_{\rm surv}(n',l')
  \label{eq:fnl}
\end{equation}
where $p_{\rm surv}(n,l)$ is the field-ionisation survival probability
($0$ or $1$ below/above the field-ionisation threshold $n_{\rm zero}$)
and $\bM$ is the radiative rate matrix.

\section{Rate Matrix}

States are indexed as $(n,l) \to i = n(n{-}1)/2 + l$, giving
$N = n_{\rm max}(n_{\rm max}{+}1)/2$ states.  The rate matrix $\bM$ is
upper-triangular (decays only go to lower~$n$):
\begin{equation}
  M_{ii} = -\frac{1}{\tau_i}\,, \qquad
  M_{ji} = \frac{B(i\!\to\!j)}{\tau_i} \quad(j < i)
\end{equation}
where $\tau_i$ is the radiative lifetime of state~$i$ and $B(i\!\to\!j)$
the branching ratio for the dipole transition $i\!\to\!j$.
Column sums of $\bM$ vanish, guaranteeing probability conservation.

\section{Method 1: Matrix Exponential}

The population vector obeys $d\bp/dt = \bM\,\bp$.  Its formal solution
$\bp(t) = \exp(\bM\,t)\,\bp(0)$ gives the exact cascade evolution for
all coupled states simultaneously.

\subsection{Implementation}

The matrix exponential is computed by scaling and squaring with a
Pad\'{e}\,[13/13] approximant~\cite{Higham2008}:
\begin{enumerate}
\item Compute the 1-norm $\|\bM\,t\|_1$ and determine a scaling
  parameter $s = \max\bigl(0,\;\lceil\log_2(\|\bM\,t\|_1/2)\rceil\bigr)$.
\item Form $\bQ = \bM\,t\,/\,2^s$.
\item Evaluate the Pad\'{e}\ numerator $\bN(\bQ)$ and denominator
  $\bD(\bQ)$ polynomials of degree~13 using Horner's rule on even/odd
  powers of $\bQ^2$.
\item Solve $\bD\,\bR = \bN$ by Newton iteration
  $\bR_{k+1} = \bR_k(2\bI - \bD\,\bR_k)$, starting from $\bR_0 = \bI$.
  Three iterations suffice ($\|\bI - \bD\|$ is small after scaling).
  Then $\bR \leftarrow \bR\,\bN$ to obtain $\bR = \bD^{-1}\bN \approx \exp(\bQ)$.
\item Square $s$~times: $\bR \leftarrow \bR^2$ to recover
  $\exp(\bM\,t)$.
\end{enumerate}
Finally, $f_{nl}$ is obtained from Eq.~(\ref{eq:fnl}) as a single
matrix--vector product $\bR^{\!\mathsf T}\,p_{\rm surv}$.
For an extended creation region ($x_1 \ne x_2$) a trapezoidal average
of $\exp(\bM\,t_1)$ and $\exp(\bM\,t_2)$ is used.

\subsection{Complexity}

Matrix multiplication dominates at $O(N^3)$ per multiply.
For $n_{\rm max} = 20$ ($N = 210$): the scaling factor is
$s \approx 15$, giving roughly $20$ matrix multiplications in total.
The $\mathcal{O}(N^3)$ cost is modest for $N \lesssim 1000$ but becomes
prohibitive for very large $n_{\rm max}$.

\section{Method 2: Recursive Cascade}

The original \texttt{hydrocal} recursive method builds $f_{nl}$ analytically, state by state,
without ever forming the full matrix \cite{Schippers2001c}.

\subsection{Core formula}

For an initial state $(n_1,l_1)$ the detection probability is
\begin{equation}
  f_{n_1 l_1} = \bigl(1 - p_{\rm d}(n_1 l_1)\bigr)\,p_{\rm surv}(n_1 l_1)
  \;+\; \sum_{n_2,l_2} B(n_1 l_1 \!\to\! n_2 l_2)\;
  \Bigl[\bigl(p_{\rm d}(n_1 l_1) - p_{\rm casc}\bigr)\,p_{\rm surv}(n_2 l_2)
  + f_{\rm casc}\Bigr]
  \label{eq:recursive}
\end{equation}
where $p_{\rm d}(nl) = 1 - e^{-t/\tau_{nl}}$ is the single-state decay
probability and $p_{\rm casc}$ is the effective cascade-chain decay
probability computed by the Lagrange-type identity:
\begin{equation}
  p_{\rm casc} = \sum_{i=0}^{K} \frac{p_{{\rm d},i}}
  {\prod_{k \ne i}\bigl(1 - \tau_k/\tau_i\bigr)}
  \label{eq:pcasc}
\end{equation}
for a chain of $K{+}1$ states with lifetimes $\tau_0, \tau_1, \ldots,
\tau_K$ and individual decay probabilities $p_{{\rm d},i}$.
The term $f_{\rm casc}$ is the recursive contribution from deeper cascade
steps, computed by calling \texttt{cascade()} with the chain extended by
one state.

The recursion depth is controlled by \texttt{ncascstep}:
\begin{itemize}
\item \texttt{ncascstep = 0}: no cascade (direct survival only).
\item \texttt{ncascstep $> 0$}: at most that many cascade steps.
\item \texttt{ncascstep $< -1$}: unlimited ($= n_{\rm max}$).
\item \texttt{ncascstep = $-1$}: selects the matrix exponential method.
\end{itemize}

\subsection{Complexity}

The recursion visits each $(n,l)$ state a limited number of times.
The total work scales roughly as $O(N \times \langle n_{\rm cascade}\rangle
\times n_{\rm neighbours})$, where $\langle n_{\rm cascade}\rangle$ is the
average chain length and $n_{\rm neighbours} \le 2$ (dipole selection
rules allow at most two $l$-values per $n_2$).  For $n_{\rm max} = 20$
this is very fast ($\sim 10^5$ operations).

\section{Comparison}

Both methods were verified against each other for $n_{\rm max} = 20$,
$q_{\rm ion} = 3$, and realistic CSRe parameters
($v_{\rm ion} = 2.0 \times 10^8$~cm/s, $L = 68$~cm).
The maximum absolute difference across all 210~states is
$\sim 3 \times 10^{-7}$, consistent with accumulated floating-point
round-off in two independent algorithms.

\begin{table}[h]
\centering
\begin{tabular}{lcc}
\toprule
 & Matrix exponential & Recursive cascade \\
\midrule
Mathematical basis & Exact ODE solution & Lagrange interpolation identity \\
Complexity & $O(N^3)$ & $O(N \times n_{\rm chain})$ \\
Matrix storage & $N \times N$ & $O(N)$ \\
Arbitrary $n_{\rm max}$ & Expensive for large $N$ & Always fast \\
Exactness & Machine precision & Machine precision \\
Truncation control & None (exact) & Via \texttt{ncascstep} \\
Extended creation & Trapezoidal average & Trapezoidal average \\
\bottomrule
\end{tabular}
\caption{Comparison of the two cascade calculation methods.}
\label{tab:comparison}
\end{table}

\subsection{Recommendations}

\begin{itemize}
\item For $n_{\rm max} \le 200$ both methods are fast.  The matrix
  exponential is simpler to implement and verify, and gives the exact
  answer in one shot.
\item The recursive method is preferred when $n_{\rm max}$ is large
  ($\gtrsim 500$), since its cost scales linearly in~$N$ rather than
  cubically.  The \texttt{ncascstep} parameter provides a natural
  convergence knob.
\item Both methods handle the trapezoidal position average for extended
  creation regions identically.
\end{itemize}

\begin{thebibliography}{1}
\bibitem{Higham2008}
  N.~J.~Higham,
  \href{https://de.scribd.com/document/401677222/Nicholas-J-Higham-Functions-of-Matrices-Theory-BookFi-org-pdf}{\emph{Functions of Matrices: Theory and Computation}}, SIAM - Society for Industrial and Applied Mathematics, Philadelpphia, USA (2008).
  
  
\bibitem{Schippers2001c}
  S. Schippers, A. Müller, G. Gwinner, J. Linkemann, A.A. Saghiri, and A. Wolf, \emph{Storage ring measurement of the C IV recombination rate coefficient}, \href{https://doi.org/10.1086/321512}{Astrophys. J. \textbf{555}, 1027 (2001)}.
\end{thebibliography}

\end{document}
